% !TEX root = ../analisi-dei-capitolati.tex
% Capitolo 1 - Introduzione

\section{Introduzione}

\subsection{Scopo del documento}
Il documento riporta l'analisi dei \glo[capitolato]{capitolati} d'appalto del
progetto didattico di Ingegneria del Software 2026/2027 e motiva la scelta
del Git Push Team (Gruppo~9): il capitolato \textbf{C2 -- NEXUM
Concierge} di Eggon S.r.l. Sono analizzati in dettaglio anche gli altri
due capitolati della rosa finale, \textbf{C1 -- AURA} (Vimar S.p.A.) e
\textbf{C0 -- AIRA} (Sanmarco Informatica S.p.A.), mentre i restanti sono
descritti più brevemente.

I termini contrassegnati dalla lettera G a pedice, come \glo{RAG}, sono
spiegati nel Glossario.

\subsection{Riferimenti}
\begin{itemize}
  \item capitolati C0--C9 dell'anno accademico 2026/2027; per C2 il
        \emph{Business Requirements Document} BRD-AI-CONCIERGE-2026,
        versione~9 del 13 ottobre 2026;
  \item verbali interni VI~2026-10-12, VI~2026-10-13 e VI~2026-10-19;
  \item verbali esterni VE~2026-10-21 (Sanmarco Informatica),
        VE~2026-10-22 (Vimar) e VE~2026-10-26 (Eggon).
\end{itemize}
